\ifdefined\standalone 
\documentclass[12pt]{article}
\usepackage{graphicx}
\usepackage{float}
\usepackage[a4paper, margin=1in]{geometry}
\usepackage[utf8]{inputenc}
\usepackage{textcomp}
\usepackage{amsmath}
\usepackage{booktabs}
\usepackage{tikz}
\usepackage{enumitem}
\usetikzlibrary{decorations.pathmorphing,patterns}
\usepackage[hidelinks]{hyperref} % <- hypertext

\begin{document}
\fi

\section{Simple TGA}
\label{sec:simple_tga}

This verification case assesses the ability of Gpyro to accurately reproduce thermogravimetric analysis (TGA) results using a single-step, first-order reaction model.
A simple non-charring material undergoing a single degradation reaction is considered.
The reaction properties are summarized in Table~\ref{tab:tga_properties}.

The simulation is performed in a zero-dimensional configuration, equivalent to a TGA experiment.
The initial sample mass is set to $m_{A,0} = 5~\mathrm{mg}$, and the temperature is increased at a constant heating rate of $\beta = 5~\mathrm{K\,min^{-1}}$, starting from an initial temperature of $T_0 = 27^\circ\mathrm{C}$.
The total simulation time is $5000~\mathrm{s}$.

The resulting mass loss rate (MLR) obtained from the Gpyro simulation is compared with the corresponding analytical solution.
The analytical formulation of this problem is derived in Appendix~\ref{ann:analytical_TGA} and is recalled here for completeness.

\begin{equation}
\mathrm{MLR}(t) =
m_{A,0}\,
Z \exp\!\left(-\frac{E}{R\,(T_0 + \beta t)}\right)
\exp\!\left(
- \int_{t_0}^{t}
Z \exp\!\left(-\frac{E}{R\,(T_0 + \beta t')}\right)\, dt'
\right)
\end{equation}

\begin{table}[H]
\centering
\caption{Reaction properties.}
\label{tab:tga_properties}
\renewcommand{\arraystretch}{1.3}
\begin{tabular}{c c c c}
\hline
Reaction &
$Z$ ($\mathrm{s^{-1}}$) &
$E$ ($\mathrm{kJ\,mol^{-1}}$) &
$n$ (-) \\
\hline
Material A $\rightarrow$ gases & $0.87 \times 10^{13}$ & 163 & 1 \\
\hline
\end{tabular}
\end{table}

The results obtained using Gpyro are compared to the analytical predictions.
Figure~\ref{fig:MLR_TGA} presents the mass loss rate as a function of time for both approaches.

\begin{figure}[H]
\centering
\includegraphics[width=0.75\linewidth]{TGA_simple.png}
\caption{Comparison between Gpyro and analytical mass loss rate predictions.}
\label{fig:MLR_TGA}
\end{figure}

\ifdefined\standalone
\end{document}
\fi




